% Midterm Section 6: Research Idea 2 (plain paragraphs without \paragraph headings; one single-column table; equations kept).
% Numbers: results/idea2_midterm/idea2_midterm_results.json (scripts/idea2_midterm.py, job 49419172)
%          session probe: results/idea2_midterm/session_probe.json
% Needs:   amsmath, booktabs, graphicx (for the figures), natbib.
% Figures: idea2_figure1.pdf, idea2_figure2.pdf next to the main tex (Overleaf root) (scripts/make_figure_idea2_midterm_v2.py)
% Bib keys: agrawal2025seamless (Seamless Interaction), ghorbani2023zeroeggs (Ghorbani et al., ZeroEGGS: Zero-shot Example-based Gesture Generation from Speech, Computer Graphics Forum 2023).

\section{Research Idea 2: Episodic Multi-Reference Conditioning}

\subsection{Initial Implementation}

Idea~2 conditions motion generation on several reference clips of the target person. Our baseline is example-based style conditioning as in ZeroEGGS~\citep{ghorbani2023zeroeggs}, which encodes one reference clip into a style vector and averages any additional clips without regard to their source recording. One clip covers one conversation, so its gestures mix the person's habits with what that conversation asked for. Clips from several recordings keep the person fixed and vary the conversation. Since Idea~2 changes only the reference side, we test that side alone before training a generator, as an identification task: an encoder must match a reference set to a held-out window of the same person, and we vary the number of clips and recordings in the set.

For a target person $p$, a reference set $R_p = \{r_1, \dots, r_K\}$ holds $K$ clips with total duration $D$ from $M \le K$ recordings, excluding the query recording. Clips from one recording start at random positions at least 10\,s apart, so four 5\,s clips are not one contiguous 20\,s segment cut in four. Each clip is cut into 5\,s blocks $b$. A block encoder $f_\theta$ maps each block to $h_b = f_\theta(b)$, and the set embedding is
\begin{equation}
  z(R_p) = \sum_{b \in B(R_p)} \alpha_b\, h_b ,
\end{equation}
with $\alpha_b = 1/|B(R_p)|$ for mean pooling. A query $q$ is a held-out 10\,s window of $p$ with embedding $e_q$ from the same encoder. Negative sets $R_{p'}$ come from other participants, have the same $(K, M, D)$ structure, and are matched to $R_p$ in voice-activity occupancy. The metric is the pairwise accuracy
\begin{equation}
  A(q) = \Pr\big[\, \lVert e_q - z(R_{p'}) \rVert > \lVert e_q - z(R_p) \rVert \,\big],
\end{equation}
averaged over matched negative sets. Chance is $A = 0.5$.

We compare three reference encoders. The frozen baseline averages the per-frame outputs of the encoders released with Seamless Interaction~\citep{agrawal2025seamless} within each block: the 64-d frame latent, the 152-d expression code with facial action units, and the 5120-d hypernetwork features. The hypernetwork features barely change within a recording and serve as an appearance control. A handcrafted encoder uses wrist speed, wrist excursion, hand height, and head angular speed in the chest frame. The learned encoder is a two-layer Transformer ($d = 128$) over upper-body SMPL-H pose at 10\,fps with temporal mean pooling. Each training episode draws two reference sets $R_p, R'_p$ per person from disjoint recordings, and the encoder minimizes
\begin{equation}
  \mathcal{L}_{\mathrm{NCE}} = -\frac{1}{P} \sum_{p=1}^{P} \log
  \frac{s_{pp}}{\sum_{p'=1}^{P} s_{pp'}},
\end{equation}
with $s_{pp'} = \exp(\langle z(R_p), z(R'_{p'})\rangle/\tau)$, symmetrized over the two sets, with $P = 48$, $\tau = 0.1$, and 4{,}000 AdamW steps. A dynamics variant subtracts each block's mean posture from the input, and a static-posture encoder uses that mean posture alone with no training.

Evaluation uses 32 participants from the first ten Naturalistic training archives, split into four participant folds. Learned encoders are trained on the other three folds plus 141 participants from later batches. Queries come from the 19 participants with at least four eligible recordings (1{,}467 windows) at a fixed 20\,s reference budget. We report participant means with 95\% participant-bootstrap intervals.

\subsection{Results and Discussion}
\label{sec:idea2_results}

\begin{table}[t]
\centering
\small
\caption{Pairwise accuracy $A$ on 19 held-out participants with a 20\,s reference budget: one 20\,s clip, four 5\,s clips of one recording, and four 5\,s clips of three recordings.}
\label{tab:idea2}
\begin{tabular}{lccc}
\toprule
Encoder & $1{\times}20$\,s & $4{\times}5$\,s & $4{\times}5$\,s \\
 & 1 rec. & 1 rec. & 3 rec. \\
\midrule
Handcrafted            & 0.60 & 0.63 & 0.64 \\
Static posture         & 0.82 & 0.84 & 0.89 \\
SI expression + FAU    & 0.92 & 0.94 & 0.96 \\
SI frame latent        & 0.93 & 0.94 & 0.95 \\
SI hypernetwork        & 0.96 & 0.96 & 0.98 \\
Learned, full pose     & 0.92 & 0.93 & 0.95 \\
Learned, dynamics only & 0.93 & 0.93 & 0.95 \\
\bottomrule
\end{tabular}
\end{table}

\begin{figure}[t]
\centering
\includegraphics[width=\columnwidth]{idea2_figure1.pdf}
\caption{Effect of contaminated references. Change in $A$ when two or four corrupted blocks are added to four clean blocks. Jittered blocks carry per-frame rotation noise, and intruder blocks come from another participant. Error bars are 95\% participant-bootstrap intervals.}
\label{fig:idea2_contam}
\end{figure}

\begin{figure}[t]
\centering
\includegraphics[width=\columnwidth]{idea2_figure2.pdf}
\caption{Accuracy per participant with one 20\,s reference from another recording, sorted by the dynamics encoder. Shaded columns mark the two participants recorded in two sessions.}
\label{fig:idea2_session}
\end{figure}

Table~\ref{tab:idea2} summarizes the results. With one 20\,s reference, the dynamics encoder reaches $A = 0.93$ and the handcrafted descriptors 0.60. Subtracting the block-mean posture does not lower accuracy (0.93 for dynamics, 0.92 for full pose), so the learned encoder identifies people from movement within a block. Posture alone reaches 0.82 with no training. The frozen appearance features score highest: the hypernetwork features reach 0.96, and 1.00 with references from the query's own recording. Same-recording references raise $A$ by 0.03 to 0.08 for every encoder, so all other results use references from other recordings.

With the number of clips fixed at four, drawing them from three recordings instead of one raises $A$ by 0.02 for the full-pose encoder and 0.01 for the dynamics encoder, with both intervals above zero. Splitting one 20\,s clip into four clips of the same recording raises $A$ by 0.013 for the full-pose encoder (95\% interval 0.006 to 0.021) and by 0.003 for the dynamics encoder (interval includes zero). For the learned encoders, additional recordings give a larger gain than splitting clips within one recording. For the handcrafted descriptors, splitting accounts for most of the gain (0.027 of 0.035), and additional recordings add little.

We found two failure cases. The first is a change of session. P0047 was recorded in two sessions, so at least half of the candidate references for each query come from the other session, and this participant has the lowest accuracy of all 19 under both the dynamics encoder (0.74) and the appearance features (0.64), as Figure~\ref{fig:idea2_session} shows. In block-mean posture, 44\% of other participants' recordings are closer to P0047's query than P0047's own recordings from the other session. The encoders capture session-specific posture as well as person-specific movement. Nearly all other participants were recorded in one session, so the results mainly reflect within-session identification.

The second failure is contaminated references. Figure~\ref{fig:idea2_contam} shows the change in $A$ as corrupted blocks are added to a clean set. With mean pooling, four blocks of another participant added to four clean blocks cost the dynamics encoder 0.06. Four noisy blocks with 0.05\,rad per-frame rotation noise cost it 0.02, because the noise looks like fast movement. The loss grows with the number of corrupted blocks, since mean pooling cannot down-weight either kind of block. The hypernetwork features lose only 0.01 under four intruders. Appearance is the more stable identity signal when the reference set is unreliable.

We then replaced mean pooling with attention pooling. A block's weight depends on its embedding and on its cosine agreement with the other blocks. Training uses contaminated sets with a penalty on the attention mass of corrupted blocks. Noisy blocks now receive almost no weight, and four of them no longer change $A$ (Figure~\ref{fig:idea2_contam}, attention gate). The cost of four intruders falls from 0.06 to 0.03 but does not vanish. When half of the set comes from another person, agreement within the set says little about identity. On clean sets the gate leaves $A$ unchanged.

\subsection{Updated Research Idea}

The results suggest different roles for appearance and motion. Appearance provides a stable identity signal, while motion captures person-specific movement but is more sensitive to corrupted references. We therefore propose appearance-guided pooling of motion embeddings from multiple recordings. Appearance determines which blocks to trust, while the final style representation contains only motion features. We hypothesize that this improves person-specific gesture generation, especially when references contain irrelevant or corrupted motion. Each reference block provides tracked motion $x_b$ and video $v_b$, a motion encoder gives $h_b = f_\theta(x_b)$, and a frozen appearance encoder gives $u_b = g(v_b)$. The block weights combine the motion agreement $c_b$ of the gate in Section~\ref{sec:idea2_results} with an appearance consistency score,
\begin{equation}
\begin{aligned}
  \alpha_b &= \frac{\exp\!\big(w^\top [\,h_b;\ c_b;\ \kappa_b\,]\big)}{\sum_{b'} \exp\!\big(w^\top [\,h_{b'};\ c_{b'};\ \kappa_{b'}\,]\big)}, \\
  \kappa_b &= \cos\!\big(u_b,\ \operatorname{median}_{b'} u_{b'}\big).
\end{aligned}
\end{equation}
Appearance enters only the weights, and $z(R_p) = \sum_b \alpha_b h_b$ is built from motion embeddings alone. Appearance consistency is most reliable when clean blocks form a majority. We evaluate its robustness under increasing intruder fractions, including the case where half of the blocks are corrupted. The objective is
\begin{equation}
\begin{aligned}
  \mathcal{L} = {} & \mathcal{L}_{\mathrm{gen}}\big(G(a, z(R_p)), x\big) + \lambda_{\mathrm{id}}\, \mathcal{L}_{\mathrm{NCE}} \\
  & + \lambda_{q} \textstyle\sum_b \alpha_b\, \mathbf{1}[b \text{ corrupted}],
\end{aligned}
\end{equation}
where $G$ is the speech-to-motion generator and $a$ the new speech. Training episodes vary the number of source recordings $M$ at a fixed duration $D$.

We also investigate cross-session robustness. Since only two participants have recordings across sessions, we first evaluate the session effect separately. If more multi-session data become available, we will test session-adversarial training while preserving participant identity.

We evaluate generated motion with an independently trained dynamics encoder, together with motion quality and speech-motion alignment metrics. At a fixed 20\,s reference budget, we compare one versus three recordings and three pooling methods: mean pooling, motion-only attention, and cross-modal attention. We also measure robustness to injected intruders and report cross-session performance separately.
